\chapter{实验设计与结果分析}\label{chap:expt}

本章通过大规模对战模拟实验，从属性覆盖度、速度线准确性和综合胜率三个维度对所提出的方法进行验证。

\section{实验设置}

\subsection{数据集构建}

本文构建了三个数据集用于实验：

\textbf{VGC使用率数据集}。从Pokémon Showdown公共数据中采集2022年1月至2025年12月的VGC规则对战记录，共计约500万场。提取每场对战中双方使用的宝可梦及其配置信息，统计各宝可梦的使用率（Usage Rate）。

\textbf{对战模拟数据集}。基于使用率排名前150的宝可梦及它们的常用配置，构建候选宝可梦池 $\mathcal{P}$。目标防御宝可梦集合 $\mathcal{D}$ 同样从这150只宝可梦中选取。

\textbf{基准队伍数据集}。收集了VGC 2022--2025赛季各大赛事的Top 8队伍配置作为基准队伍，共计96支队伍。

\subsection{评价指标}

本文采用以下指标评估队伍优化效果：

\begin{enumerate}
    \item \textbf{属性覆盖度（Coverage Rate）}：队伍对目标宝可梦集合的克制度比例，见公式~\eqref{eq:coverage}；
    \item \textbf{速度配合度（Speed Synergy）}：队伍速度线阶梯分布的信息熵，熵值越高表示速度分布越合理；
    \item \textbf{综合胜率（Win Rate）}：在Pokémon Showdown对战模拟器中，队伍与基准队伍进行100场对战的胜率。
\end{enumerate}

\subsection{对比方法}

本文将所提出的GLSH+NSGA-II方法（记为\textbf{GLSH-NSGA}）与以下方法进行对比：

\begin{enumerate}
    \item \textbf{Random}：随机选取6只宝可梦并随机配置技能和努力值；
    \item \textbf{Usage}：选取使用率最高的6只宝可梦及其热门配置；
    \item \textbf{Greedy}：仅使用贪心策略选取宝可梦\citep{chu2004tushu}；
    \item \textbf{GA-only}：使用标准单目标遗传算法\citep{walls2013drought}，仅优化属性覆盖度。
\end{enumerate}

\section{属性覆盖度实验}

\subsection{覆盖度对比}

表~\ref{tab:coverage-results}展示了不同方法在属性覆盖度上的表现。

\begin{table}[!htbp]
    \bicaption{属性覆盖度对比}{Type coverage comparison of different methods}
    \label{tab:coverage-results}
    \centering
    \footnotesize
    \begin{tabular}{lcccc}
        \hline
        方法 & 覆盖度(\%) & 平均克制倍率 & 最大覆盖缺口 \\
        \hline
        Random & 62.3 & 0.87 & 56 \\
        Usage & 78.9 & 1.12 & 32 \\
        Greedy & 84.5 & 1.23 & 23 \\
        GA-only & 87.1 & 1.28 & 19 \\
        \textbf{GLSH-NSGA} & \textbf{91.3} & \textbf{1.41} & \textbf{13} \\
        \hline
    \end{tabular}
\end{table}

实验结果表明，GLSH-NSGA在属性覆盖度上达到91.3\%，相比Usage方法提升12.4个百分点，相比GA-only方法提升4.2个百分点。最大覆盖缺口从56（Random）降低至13（GLSH-NSGA），表明队伍对主流防御宝可梦的克制盲区显著减小。

\subsection{双属性组合效应分析}

图~\ref{fig:dual-type}展示了考虑双属性组合前后覆盖度的差异。引入双属性组合效应因子后，覆盖度评估更为精确，避免了因忽视双属性抵抗而导致的覆盖度虚高问题。

\begin{figure}[!htbp]
    \centering
    \includegraphics[width=0.50\textwidth]{pokemon/mega_blastoise}
    \bicaption{单属性与双属性评估的覆盖度差异}{Coverage difference between single-type and dual-type evaluation}
    \label{fig:dual-type}
\end{figure}

\section{速度线预测实验}

\subsection{预测准确性}

本文使用留一赛季交叉验证（Leave-One-Season-Out）评估速度线预测模型的准确性。以VGC 2022--2024赛季数据训练，预测2025赛季关键速度门槛，并与实际速度线分布进行对比。

实验结果显示，贝叶斯优化模型预测的速度线阈值与实际关键门槛的平均绝对误差（MAE）为3.2点速度值，优于基于简单统计的方法（MAE = 7.8点速度值）。

\subsection{速度配合度分析}

表~\ref{tab:speed-synergy}对比了不同方法生成队伍的速度配合度。

\begin{table}[!htbp]
    \bicaption{速度配合度对比}{Speed synergy comparison}
    \label{tab:speed-synergy}
    \centering
    \footnotesize
    \begin{tabular}{lcc}
        \hline
        方法 & 速度配合度 & 信息熵 \\
        \hline
        Random & 0.32 & 1.21 \\
        Usage & 0.58 & 1.89 \\
        Greedy & 0.61 & 1.94 \\
        GA-only & 0.55 & 1.76 \\
        \textbf{GLSH-NSGA} & \textbf{0.78} & \textbf{2.34} \\
        \hline
    \end{tabular}
\end{table}

GLSH-NSGA在多目标优化框架中显式考虑了速度配合度，生成队伍的速度线分布更加合理，避免了单一速度段宝可梦堆叠的问题。

\section{综合胜率实验}

\subsection{对战模拟设置}

在Pokémon Showdown平台上，使用GLSH-NSGA生成的最优队伍与基准队伍进行对战测试。每种对比进行100场对战，消除随机因素（如命中率、会心率）带来的方差。

\subsection{胜率对比}

表~\ref{tab:winrate}展示了各方法生成队伍的胜率对比。

\begin{table}[!htbp]
    \bicaption{综合胜率对比}{Win rate comparison (100 battles)}
    \label{tab:winrate}
    \centering
    \footnotesize
    \begin{tabular}{lcccc}
        \hline
        方法 & 胜 & 负 & 胜率(\%) & 平均剩余精灵数 \\
        \hline
        Random & 23 & 77 & 23.0 & 1.2 \\
        Usage & 47 & 53 & 47.0 & 2.3 \\
        Greedy & 52 & 48 & 52.0 & 2.5 \\
        GA-only & 55 & 45 & 55.0 & 2.7 \\
        \textbf{GLSH-NSGA} & \textbf{68} & \textbf{32} & \textbf{68.0} & \textbf{3.4} \\
        \hline
    \end{tabular}
\end{table}

GLSH-NSGA的综合胜率达到68.0\%，显著优于其他方法。平均剩余精灵数达到3.4只，表明队伍在对战中具有较好的容错性。

\section{消融实验}

为分析GLSH-NSGA各组件（贪心构建、局部搜索、速度优化、攻防平衡）的贡献，进行了消融实验，结果如表~\ref{tab:ablation}所示。

\begin{table}[!htbp]
    \bicaption{消融实验结果}{Ablation study results}
    \label{tab:ablation}
    \centering
    \footnotesize
    \begin{tabular}{lcccc}
        \hline
        配置 & 覆盖度(\%) & 速度配合度 & 胜率(\%) \\
        \hline
        GLSH-NSGA（完整） & 91.3 & 0.78 & 68.0 \\
        $-$ 局部搜索 & 87.6 & 0.77 & 62.5 \\
        $-$ 速度优化 & 90.8 & 0.61 & 61.0 \\
        $-$ 攻防平衡 & 90.5 & 0.75 & 64.2 \\
        $-$ 贪心构建 & 85.0 & 0.72 & 58.8 \\
        \hline
    \end{tabular}
\end{table}

消融实验表明：局部搜索对覆盖度贡献最大（$-$3.7\%），速度优化对胜率贡献最大（$-$7.0\%），验证了各组件在整体框架中的必要性和有效性。

\section{本章小结}

本章通过属性覆盖度实验、速度线预测实验和综合胜率实验验证了GLSH-NSGA方法的有效性。实验结果表明，该方法在属性覆盖度（91.3\%）、速度配合度（0.78）和综合胜率（68.0\%）三个指标上均显著优于对比方法。消融实验进一步验证了各算法组件的独立贡献。
